\documentclass[usepdftitle=false]{beamer}

\usepackage{amssymb}
\usepackage{mathrsfs}
\usepackage{stmaryrd}
\usepackage{tikz-cd}
\usepackage{ulem}
\usepackage{algpseudocode}

\usepackage[dvipsnames]{xcolor}

\newcommand{\CC}{\mathbb{C}}
\newcommand{\FF}{\mathbb{F}}
\newcommand{\NN}{\mathbb{N}}
\newcommand{\QQ}{\mathbb{Q}}
\newcommand{\RR}{\mathbb{R}}
\newcommand{\ZZ}{\mathbb{Z}}
\newcommand{\id}{\mathrm{id}}
\newcommand{\op}[1]{#1^{\text{op}}}
\newcommand{\powerset}[1]{\mathcal{P}(#1)}

\newcommand{\Set}{\mathbf{Set}}

\newcommand{\true}{\mathrm{true}}
\newcommand{\false}{\mathrm{false}}

\newcommand{\varset}[0]{\mathbb{V}}
\newcommand{\valset}[0]{\mathbb{I}}
\newcommand{\stateset}[0]{\mathcal{E}}
\newcommand{\concdomain}[0]{\mathcal{D}}
\newcommand{\statesetb}[0]{\hat{\mathcal{E}}}
\newcommand{\concdomainb}[0]{\hat{\mathcal{D}}}
\newcommand{\statset}[0]{\mathbb{S}}

\newcommand{\expression}[1]{\textcolor{red}{\mathbb{E}\llbracket} #1 \textcolor{red}{\rrbracket}}
\newcommand{\expressionc}[1]{\textcolor{red}{\mathbb{E'}\llbracket} #1 \textcolor{red}{\rrbracket}}
\newcommand{\cexpression}[1]{\textcolor{red}{\mathbb{C}\llbracket} #1 \textcolor{red}{\rrbracket}}
\newcommand{\cexpressionc}[1]{\textcolor{red}{\mathbb{C'}\llbracket} #1 \textcolor{red}{\rrbracket}}
\newcommand{\statement}[1]{\textcolor{red}{\mathbb{S}\llbracket} #1 \textcolor{red}{\rrbracket}}
\newcommand{\statementb}[1]{\textcolor{red}{\hat{\mathbb{S}}\llbracket} #1 \textcolor{red}{\rrbracket}}
\newcommand{\statementc}[1]{\textcolor{red}{\mathbb{S'}\llbracket} #1 \textcolor{red}{\rrbracket}}
\newcommand{\absexpression}[1]{\textcolor{red}{\mathbb{E}^\sharp\llbracket} #1 \textcolor{red}{\rrbracket}}
\newcommand{\abscexpression}[1]{\textcolor{red}{\mathbb{C}^\sharp\llbracket} #1 \textcolor{red}{\rrbracket}}
\newcommand{\absstatement}[1]{\textcolor{red}{\mathbb{S}^\sharp\llbracket} #1 \textcolor{red}{\rrbracket}}

\newcommand{\ite}[3]{\mathbf{if} \, #1 \, \mathbf{then} \, #2 \, \mathbf{else} \, #3 \, \mathbf{end\,if}}
\newcommand{\while}[2]{\mathbf{while} \, #1 \, \mathbf{do} \, #2 \, \mathbf{end\,while}}

\DeclareMathOperator{\lfp}{lfp}

\usepackage{tikz}
\usetikzlibrary{shapes.geometric, calc}

%\renewcommand\labelenumii{(\roman{enumii})}

%\theoremstyle{definition}
%\newtheorem{theorem}{Theorem}[section]
\newtheorem{proposition}[theorem]{Proposition}
%\newtheorem{definition}[theorem]{Definition}
%\newtheorem{lemma}[theorem]{Lemma}
%\newtheorem{corollary}[theorem]{Corollary}
%\newtheorem{example}{Example}[theorem]

\usepackage[many]{tcolorbox}
\tcbuselibrary{skins}

\definecolor{enby1}{RGB}{255,244,51}
\definecolor{enby2}{RGB}{255,248,231}
\definecolor{enby3}{RGB}{155,89,208}
\definecolor{enby4}{RGB}{45,45,45}
\tcbset{enby/.style={enhanced,hbox,colback=enby1,
underlay={\begin{tcbclipinterior}
\fill[enby2] let \p1=($(interior.north)-(interior.south)$) in
([yshift={2*\y1/4}]interior.south west) rectangle 
([yshift={1*\y1/4}]interior.south east);
\fill[enby3] let \p1=($(interior.north)-(interior.south)$) in
([yshift={3*\y1/4}]interior.south west) rectangle 
([yshift={2*\y1/4}]interior.south east);
\fill[enby4] let \p1=($(interior.north)-(interior.south)$) in
([yshift={4*\y1/4}]interior.south west) rectangle 
([yshift={3*\y1/4}]interior.south east);
\end{tcbclipinterior}}},
tight/.style={on
line,boxsep=1pt,left=#1,right=#1,top=#1,bottom=#1,arc=0pt,boxrule=0pt},
tight/.default=1pt}

\mode<presentation>
{
    %\usetheme{Warsaw}
}

\title{Using AI to improve your code}
\author{Eris :3 \begin{tcolorbox}[tight, enby](they/them)\end{tcolorbox}}
\date{\today}

\begin{document}

\begin{frame}
    \titlepage
\end{frame}

\begin{frame}
    \frametitle{I FUCKING HATE AI!!!}

    \uncover<2->{
    Please imagine a video of Sam Altman being blown up by a missile.
    }

    \uncover<3->{
    These are my thoughts on ai :)
    }
    \uncover<4->{
    Artificial intelligence will not be mentioned again for the rest of this talk.
    }
\end{frame}

\title{Using \sout{AI} Abstract Interpretation to improve your code}

\begin{frame}
    \titlepage
\end{frame}

\title{How other people use Abstract Interpretation to improve your code when you run it}

\begin{frame}
    \titlepage
\end{frame}

\begin{frame}
    \frametitle{Optimising compilers}

    \begin{center}
        \begin{tikzpicture}
            \node (s) at (0, 0) {Source code};
            \node (m) at (6, 0) {Machine code};
            \draw[->] (s) -- (m);
            \node (c) at (3, 0) [fill=white, draw=black] {Compiler};
            \uncover<2->{
            \draw[->] (c) edge [loop above] node[above] {Optimisations} (c);
            }
        \end{tikzpicture}
    \end{center}
\end{frame}

\begin{frame}
    \frametitle{Fibonacci program}

    \only<1-1>{
    \begin{algorithmic}
        \State \(F \gets \underbrace{[0, \dots, 0]}_{\text{100 of them}}\)
            \If{\(\lnot(0 \leq 0 \leq 99 \land 0 \leq 1 \leq 99)\)}
                \State \textbf{explode} (out of bounds)
            \EndIf
        \State \(F[0] \gets 0\)
        \State \(F[1] \gets 1\)
        \State \(i \gets 2\)
        \While{\(i < 100\)}
            \If{\(\lnot(0 \leq i-2 \leq 99 \land 0 \leq i-1 \leq 99 \land 0 \leq i \leq 99)\)}
                \State \textbf{explode} (out of bounds)
            \EndIf
            \State \(F[i] \gets F[i-1] + F[i-2]\)
            \State \(i \gets i + 1\)
        \EndWhile
        \\\noindent\Return \(F\)
    \end{algorithmic}
    }

    \only<2-2>{
    \begin{algorithmic}
        \State \(F \gets \underbrace{[0, \dots, 0]}_{\text{100 of them}}\)
            \If{\(\false\)}
                \State \textbf{explode} (out of bounds)
            \EndIf
        \State \(F[0] \gets 0\)
        \State \(F[1] \gets 1\)
        \State \(i \gets 2\)
        \While{\(i < 100\)}
            \If{\(\lnot(0 \leq i-2 \leq 99 \land 0 \leq i-1 \leq 99 \land 0 \leq i \leq 99)\)}
                \State \textbf{explode} (out of bounds)
            \EndIf
            \State \(F[i] \gets F[i-1] + F[i-2]\)
            \State \(i \gets i + 1\)
        \EndWhile
        \\\noindent\Return \(F\)
    \end{algorithmic}
    }

    \only<3-3>{
    \begin{algorithmic}
        \State \(F \gets \underbrace{[0, \dots, 0]}_{\text{100 of them}}\)
        \State \(F[0] \gets 0\)
        \State \(F[1] \gets 1\)
        \State \(i \gets 2\)
        \While{\(i < 100\)}
            \If{\(\lnot(0 \leq i-2 \leq 99 \land 0 \leq i-1 \leq 99 \land 0 \leq i \leq 99)\)}
                \State \textbf{explode} (out of bounds)
            \EndIf
            \State \(F[i] \gets F[i-1] + F[i-2]\)
            \State \(i \gets i + 1\)
        \EndWhile
        \\\noindent\Return \(F\)
    \end{algorithmic}
    }

    \only<4->{
    \begin{algorithmic}
        \State \(F \gets \underbrace{[0, \dots, 0]}_{\text{100 of them}}\)
        \State \(F[0] \gets 0\)
        \State \(F[1] \gets 1\)
        \State \(i \gets 2\)
        \While{\(i < 100\)}
            \If{\(i < 2 \lor 99 < i\)}
                \State \textbf{explode} (out of bounds)
            \EndIf
            \State \(F[i] \gets F[i-1] + F[i-2]\)
            \State \(i \gets i + 1\)
        \EndWhile
        \\\noindent\Return \(F\)
    \end{algorithmic}

    \uncover<5->{
        The bounds check is redundant though!
    }
    }
\end{frame}

\begin{frame}
    \frametitle{Syntax}
    We define expressions, conditions and statements.

    \uncover<2->{
    Expressions are things that evaluate to numbers.
    \begin{itemize}
        \item<3->
            A number \(n\) is an expression.
        \item<4->
            A variable (integer variable) \(v\) is an expression.
        \item<5->
            Index an array by an expression \(a[e]\).
        \item<6->
            We can write down the sum of two expressions \(e_1 + e_2\).
        \item<7->
            Difference \(e_1 - e_2\).
        \item<8->
            Product \(e_1 \times e_2\).
    \end{itemize}
    }
    \[
        \begin{array}{rcl}
            s & := & v \gets e
                \; \mid \; s; s
                \; \mid \; \ite{c}{s}{s} \\
            & & \; \mid \; \while{c}{s}
                \; \mid \; \mathbf{explode}
                \; \mid \; \mathbf{skip}
        \end{array}
    \]
\end{frame}

\begin{frame}
    \frametitle{Syntax}
    Conditions are what guard ifs and whiles.
    \begin{itemize}
        \item<2->
            \(\true\) and \(\false\) are conditions
        \item<3->
            Comparisons of expressions like \(e_1 < e_2\) or \(e_1 = e_2\).
        \item<4->
            Conjunctions \(c_1 \land c_2\)
        \item<5->
            Disjunctions \(c_1 \lor c_2\)
    \end{itemize}
    \uncover<6->{
        Negation is dealt with by applying de-morgan's law and negating "atoms".
        \[
            \lnot(x \leq y \land y \leq z) = x > y \lor y > z
        \]
    }
\end{frame}

\begin{frame}
    \frametitle{Syntax}
    Statements are the lines of code.
    \begin{itemize}
        \item<2->
            Variable assignment by expressions \(v \gets e\)
        \item<3->
            Array updates, indexed by expressions \(a[e_1] \gets e_2\)
        \item<4->
            The composite of two statements \(s_1; s_2\)
        \item<5->
            If statements, with bodies statements \(\ite{c}{s_1}{s_2}\)
        \item<6->
            While loops, with bodies statements \(\while{c}{s}\)
    \end{itemize}
\end{frame}

\begin{frame}
    \frametitle{Semantics}
    There is a set \(\stateset\) of states (assignments of all variables to values) and (letting \(\concdomain = \mathcal{P}(\stateset)\)) an operator \(\statement{s} : \concdomain \to \concdomain\) for each statement \(s\).
    This operator describes ``what the statements do''.

    \uncover<2->{
        For example,
        \[
            \statement{x \gets 5}(\stateset) = \{\sigma \in \stateset \mid \sigma(x) = 5\}.
        \]
    }

    \uncover<3->{
    It's quite yucky to write out everything in full with arrays included.
    Instead, we will use order theory to ``approximate'' the real semantics.
    }
\end{frame}

\begin{frame}
    \frametitle{Order theory 1}
    For any set \(S\), \((\mathcal{P}(S), \subseteq)\) is a poset.

    \uncover<2->{
        \begin{center}
            \begin{tikzpicture}
                \node (top) at (0,2) {$\{x,y,z\}$};
                \node (xy) at (-1,1) {$\{x,y\}$};
                \node (xz) at (0,1) {$\{x,z\}$};
                \node (yz) at (1,1) {$\{y,z\}$};
                \node (x) at (-1,0) {$\{x\}$};
                \node (y) at (0,0) {$\{y\}$};
                \node (z) at (1,0) {$\{z\}$};
                \node (bot) at (0,-1) {$\emptyset$};
                \draw (bot) -- (x) -- (xy) -- (top);
                \draw (bot) -- (y) -- (yz) -- (top);
                \draw (bot) -- (z) -- (xz) -- (top);
                \draw (x) -- (xz);
                \draw (y) -- (xy);
                \draw (z) -- (yz);
            \end{tikzpicture}
        \end{center}
    }
\end{frame}

\begin{frame}
    \frametitle{Order theory 1}
    \begin{definition}
        A map \(f : P \to Q\) between posets is monotone if whenever \(x \sqsubseteq_P y\), \(f(x) \sqsubseteq_Q f(y)\).
    \end{definition}
    \uncover<2->{
    \begin{proposition}
        The identity map is monotone. Composites of monotone maps are monotone.
    \end{proposition}
    }
\end{frame}

\begin{frame}
    \frametitle{Approximating semantics}
    Instead of defining the full semantics, define an approximation set of states \(\statesetb\) and (letting \(\concdomainb = \mathcal{P}(\statesetb)\)) define an approximation operator \(\statementb{s} : \concdomainb \to \concdomainb\) for each statement \(s\).

    \uncover<2->{
    Need a monotone map \(\hat{\gamma} : \concdomainb \to \concdomain\) such that for any set of states \(S\), \(\statement{s} (\gamma(S)) \subseteq \gamma (\statementb{s}(S))\).
    Intuitively, every state described in the ``real world'' is captured in the ``approximate world''.
    }
    
    \uncover<3->{
    If you prefer, we need \(\hat{\gamma}\) a lax natural transformation viewing semantics opeators as functors from the one object category of statements under sequential composition to the poset-enriched category of posets!
    \[
    \begin{tikzcd}[ampersand replacement=\&]
        \concdomainb \arrow[r, "\statementb{s}"] \arrow[d, "\gamma", swap] \& \concdomainb \arrow[d, "\gamma"] \\
        \concdomain \arrow[r, "\statement{s}", swap] \arrow[ur, "\sqsubseteq" sloped, phantom] \& \concdomain
    \end{tikzcd}
    \]
    }
\end{frame}

\begin{frame}
    \frametitle{Approximating semantics}
    Define \(\statesetb = \ZZ^V\) (where \(V\) is the set of \textit{integer} variables only) and \(\concdomainb = \mathcal{P}(\statesetb)\).
    Then define \(\gamma : \concdomainb \to \concdomain\), to give all possible assignments of arrays.

    Then we'll define \(\expression{e} : \statesetb \to \powerset{\ZZ}\), \(\cexpression{c} : \concdomainb \to \concdomainb\) and \(\statementb{s} : \concdomainb \to \concdomainb\) inductively.

    \(\expression{e}\) tells us what expressions can evaluate to and \(\cexpression{c}\) tells us which statements pass through a guard.
\end{frame}

\begin{frame}
    \frametitle{Approximating semantics}
    \[
        \begin{array}{rcl}
            \expression{n}(\sigma) &=& \{n\} \\
            \expression{v}(\sigma) &=& \{\sigma(v)\} \\
            \expression{e_1 + e_2}(\sigma) &=& \{x + y \mid x \in \expression{e_1}(\sigma), y \in \expression{e_2}(\sigma)\} \\
            \expression{e_1 - e_2}(\sigma) &=& \{x - y \mid x \in \expression{e_1}(\sigma), y \in \expression{e_2}(\sigma)\} \\
            \expression{e_1 \times e_2}(\sigma) &=& \{x \times y \mid x \in \expression{e_1}(\sigma), y \in \expression{e_2}(\sigma)\} \\
            \expression{a[i]}(\sigma) &=& \ZZ \\
        \end{array}
    \]
    \[
        \begin{array}{rcl}
            \cexpression{\true}(S) &=& S \\
            \cexpression{\false}(S) &=& \emptyset \\
            \cexpression{e_1 \leq e_2}(S) &=& \{\sigma \in S \mid \exists x \in \expression{e_1}, y \in \expression{e_2} : x \leq y\} \\
            \cexpression{e_1 = e_2}(S) &=& \{\sigma \in S \mid \exists x \in \expression{e_1}, y \in \expression{e_2} : x = y\} \\
            \cexpression{c_1 \land c_2}(S) &=& \cexpression{c_1}(S) \cap \cexpression{c_2}(S) \\
            \cexpression{c_1 \lor c_2}(S) &=& \cexpression{c_1}(S) \cup \cexpression{c_2}(S) \\
        \end{array}
    \]
\end{frame}

\begin{frame}
    \frametitle{Approximating semantics}
    \[
        \begin{array}{l}
            \statementb{s_1; s_2}(S) = \statementb{s_2}(\statementb{s_1}(S)) \\
            \statementb{v \gets e}(S) = \{\sigma[v \mapsto x] \mid \sigma \in S, x \in\expression{e}(\sigma)\} \\
            \statementb{a[i] \gets e}(S) = S \\
            \statementb{\ite{c}{s_1}{s_2}}(S) =\\ \qquad\qquad\qquad \statementb{s_1}(\cexpression{c}(S)) \cup  \statementb{s_1}(\cexpression{\lnot c}(S)) \\
            \uncover<2->{
            \statementb{\while{c}{s}}(S) = \text{wait what} \\
            }
        \end{array}
    \]
    \noindent\uncover<3->{
    We need some sort of limit of repeatedly applying \(s\)!
    Alternatively, we could find some minimal \(X\) such that \(X = X \cup \statementb{s}(\cexpression{c}(X))\).
    }
\end{frame}

\begin{frame}
    \frametitle{Order theory 2}
    Unfortunately domain theorists use \(\omega\) to mean \(\NN\).

    \uncover<2->{
    \begin{definition}
        An \(\omega\)-chain is a sequence \((x_i)_{i \in \omega}\) such that if \(i \sqsubseteq j\) then \(x_i \sqsubseteq x_j\).
    \end{definition}
    }
    \uncover<3->{
    \begin{definition}
        An \(\omega\) complete partial order, or \(\omega\)-CPO is a poset \((C, \sqsubseteq)\) such that all \(\omega\)-chains \((x_i)\) admit a least upper bound \(\bigsqcup_i x_i\) (so \(x_j \sqsubseteq \bigsqcup_i x_i\) for all \(j\) and is minimal among such).
        We may call these least upper bounds \textit{joins}.
        An \(\omega\)-CPO is \textit{pointed} if it has a least element \(\bot\).
    \end{definition}
    }
\end{frame}

\begin{frame}
    \frametitle{Order theory 2}
    \begin{example}
        For any set \(X\), the poset \((\mathcal{P}(X), \subseteq)\) is an \(\omega\)-cpo!
        Joins are given by unions of sets, \(\bigsqcup_i X_i = \bigcup_i X_i\).
        Moreover, it is pointed by \(\emptyset\)!
    \end{example}
    We can think of joins as generalisations of unions, but we will only see instances of non-union joins later.
\end{frame}

\begin{frame}
    \frametitle{Order theory 2}
    \begin{definition}
        A monotone map \(f : C \to D\) of \(\omega\)-cpos is \(\omega\)-continuous if for any \(\omega\)-chain \((x_i)_{i \in \omega}\) in \(C\), \(f(\bigsqcup_i x_i) = \bigsqcup_i f(x_i) \).
    \end{definition}
    \begin{theorem}[Kleene fixed-point theorem for \(\omega\)-cpos]
        If \(C\) is a pointed \(\omega\)-cpo and \(f : C \to C\) is an \(\omega\)-continuous map, then a fixed point of \(f\) exists, and the least fixed point is
        \[
            \lfp(f) = \bigsqcup_i f^i(\bot).
        \]
    \end{theorem}
\end{frame}

\begin{frame}
    \frametitle{Order theory 2}
    \begin{proof}
        Notice that \(f^i\) is monotone since it is a composite of monotone maps.
        Applied to \(\bot \sqsubseteq f(\bot)\), we get \(f^i(\bot) \sqsubseteq f^{i+1}(\bot)\), so \(f^i(\bot)\) does form an \(\omega\)-chain!

        \uncover<2->{
        Next, our computation does give a fixed point, since
        \[
            f\left(\bigsqcup_i f^i(\bot)\right) = \bigsqcup_i f^{i+1}(\bot) = \bigsqcup_i f^i(\bot).
        \]
        }

        \uncover<3->{
        Finally, if \(x\) is another fixed-point of \(f\), then applying monotonicity to \(\bot \sqsubseteq x\) we get
        \(f^i(\bot) \sqsubseteq f^i(x) = x\), so \(\sqsubseteq_i f^i(\bot) \sqsubseteq x\).
        }
    \end{proof}
\end{frame}

\begin{frame}
    \frametitle{While loops}
    And now we can define while loop semantics!
    \uncover<2->{
    \[
        \begin{array}{l}
            \statementb{\while{c}{s}}(S) =\\\qquad\qquad\qquad \cexpression{\lnot c}(\lfp (X \mapsto X \cup \statementb{s}(\cexpression{c}(X))))
        \end{array}
    \]
    We've finally finished defining the language ...
    }

    \noindent\uncover<3->{
    (There are questions about how we can know the inner expression is continuous if we haven't fully defined it but shhh this talk isn't about induction)
    }
\end{frame}

\begin{frame}
    \frametitle{Abstract interpretation}
    The semantics we just described is very non-computable!
    Abstract interpretation is about approximating our ``concrete'' semantics so that it becomes computable.
    
    \uncover<2->{
    That is, we want a poset \((A, \sqsubseteq)\) with elements of \(A\) easy to represent on a computer, and some computable semantics operator \(\absstatement{s} : A \to A\) for each statement \(s\), that's an approximation of our semantics by a map \(\gamma : A \to \concdomainb\).
    }
    
    \uncover<3->{
    One approach is to store the sets of all values a variable could have (\(\mathcal{D}' = \mathcal{P}(\ZZ)^V\)).
    We can then approximate \(\mathcal{P}(\ZZ)\) by \(A\), and then get some computable set \(A^V\)!
    }
\end{frame}

\begin{frame}
    \frametitle{Value domains}
    Examples of approximations of \(\mathcal{P}(\ZZ)\):
    \begin{itemize}
        \item<2->
            \(S = \begin{tikzpicture}[scale=.7, baseline=($(zer)!0.15!(bot)$)]
                \node (top) at (0,1) {$\top$};
                \node (neg) at (-1,0) {$-$};
                \node (zer) at (0,0) {$0$};
                \node (pos) at (1,0) {$+$};
                \node (bot) at (0,-1) {$\bot$};
                \draw (bot) -- (neg) -- (top);
                \draw (bot) -- (zer) -- (top);
                \draw (bot) -- (pos) -- (top);
            \end{tikzpicture}\).

            \(\gamma(+) = \{n \in \ZZ \mid n > 0\}\) etc.
        \item<3->
            \(A = \{[a, b] \mid a, b \in \ZZ \cup \{-\infty, \infty\}\}\) ordered by subsets.

            \(\gamma([a, b]) = \{n \in \ZZ \mid a \leq n \leq b\}\).
    \end{itemize}
\end{frame}

\begin{frame}
    \frametitle{Example again}
    \begin{algorithmic}
        \State \(F \gets \underbrace{[0, \dots, 0]}_{\text{100 of them}}\)
        \State \(F[0] \gets 0\)
        \State \(F[1] \gets 1\)
        \State \(i \gets 2\)
        \While{\(i < 100\)}
            \If{\(2 > i \lor i > 99\)}
                \State \textbf{explode} (out of bounds)
            \EndIf
            \State \(F[i] \gets F[i-1] + F[i-2]\)
            \State \(i \gets i + 1\)
        \EndWhile
        \\\noindent\Return \(F\)
    \end{algorithmic}
    At the second bounds check, \(i \in [2, 99]\)!!
\end{frame}

\begin{frame}
    \frametitle{Value set semantics}
    We define just \(\expressionc{e} : \mathcal{P}(\ZZ)^V \to \mathcal{P}(\ZZ)\) and leave the rest for later.

    \uncover<2->{
    \[
        \begin{array}{rcl}
            \expressionc{n}(\sigma) &=& \{n\} \\
            \expressionc{v}(\sigma) &=& \sigma(v) \\
            \expressionc{e_1 + e_2}(\sigma) &=& \{x + y \mid x \in \expressionc{e_1}(\sigma), y \in \expressionc{e_2}(\sigma)\} \\
            \expressionc{e_1 - e_2}(\sigma) &=& \{x - y \mid x \in \expressionc{e_1}(\sigma), y \in \expressionc{e_2}(\sigma)\} \\
            \expressionc{e_1 \times e_2}(\sigma) &=& \{x \times y \mid x \in \expressionc{e_1}(\sigma), y \in \expressionc{e_2}(\sigma)\} \\
            \expressionc{a[i]}(\sigma) &=& \ZZ
        \end{array}
    \]
    }
    \uncover<3->{
    We want to approximate \(\expressionc{e}\) with intervals instead of sets.
    }
\end{frame}

\begin{frame}
    \frametitle{Order theory 3}
    oops this is out of place but needs to go somewhere
    \uncover<2->{
    \begin{definition}[Non-binary double feature AND lattice double feature?!]
        A poset \((P, \sqsubseteq)\) is a join semi-lattice if every pair of elements \(x, y\) has a least upper bound (join) written \(x \sqcup y\).
    \end{definition}
    }
    \uncover<3->{
    \begin{example}
        As usual \((\mathcal{P}(X), \sqsubseteq)\) is a join semi-lattice, where joins are unions!
    \end{example}
    }
\end{frame}

\begin{frame}
    \frametitle{Order theory 3}
    \begin{definition}
        If \((C,\sqsubseteq)\) and \((A, \sqsubseteq)\) are posets, and \(\alpha : C \rightleftarrows A : \gamma\) are monotone maps, then they form a Galois connection if
        \(\alpha(c) \sqsubseteq a\) iff \(c \sqsubseteq \gamma(a)\).
        We call \(\alpha\) the left adjoint and \(\gamma\) the right adjoint.
    \end{definition}
    \uncover<2->{
    \begin{example}
        The map \(\gamma([a, b]) = \{n \in \ZZ \mid a \leq n \leq b\}\) has left adjoint \(\alpha(S) = [\min(S), \max(S)]\)!
    \end{example}
    }
\end{frame}

\begin{frame}
    \frametitle{Order theory 3}
    Galois embeddings let us optimally approximate concrete structure in the abstract structure.
    \uncover<2->{
    \begin{proposition}
        In a Galois connection \(\alpha : C \rightleftarrows A : \gamma\), if \(C\) is a join semi-lattice, then \(A\) is a join semi-lattice with join \(a \sqcup b = \alpha(\gamma(a) \sqcup \gamma(b))\).
    \end{proposition}
    \begin{proposition}
        Any operator \(f : C \to C\) has its ``best abstraction'' on \(A\) (won't be defined sorry) by \(\alpha \circ f \circ \gamma\).
    \end{proposition}
    }
    \uncover<3->{
    (I'm too dead to write the proofs for this... do them as exercises!)
    }

    \uncover<4->{
    As a general idea, apply \(\gamma\), then some operation, then \(\alpha\).
    This gives us approximation operators for free yay!
    }
\end{frame}

\begin{frame}
    \frametitle{Computing Interval semantics}
    Let's write out what \(\absexpression{e_1 + e_2} : I^V \to I\) might be.
    \uncover<2->{
    Let \([a, b] = \absexpression{e_1}(\sigma^\sharp), [c, d] = \absexpression{e_2}(\sigma^\sharp)\).
    \begin{align*}
        \absexpression{e_1 + e_2}(\sigma^\sharp)
        &= \alpha(\expressionc{e_1 + e_2}(\gamma(\sigma^\sharp))) \\
        &\sqsubseteq \alpha(\{x + y \mid x \in [a, b], y \in [c, d]\}) \\
        &= [\min(\{x + y \mid x \in [a, b], y \in [c, d]\}), \max(\cdots)] \\
        &= [a + c, b + d] \\
    \end{align*}
    }
    \uncover<3->{
    wow it's just interval arithmetic!
    }
\end{frame}

\begin{frame}
    \frametitle{Remaining interval semantics}
    We still have to write out \(\abscexpression{c}\) and \(\absstatement{s}\) ourselves.
    \uncover<2->{
    \[
        \begin{array}{rcl}
            \abscexpression{v < n}(\sigma^\sharp) &=& \sigma[v \mapsto \sigma(v) \sqcap [-\infty, n-1]] \\
            \abscexpression{v = n}(\sigma^\sharp) &=& \sigma[v \mapsto \sigma(v) \sqcap [n, n]] \text{ etc} \\
            \abscexpression{c}(\sigma^\sharp) &=& \sigma^\sharp \\
        \end{array}
    \]
    }
    \uncover<3->{
    \[
        \begin{array}{l}
            \absstatement{v \gets e}(S) = \sigma[v \mapsto \absexpression{e}(\sigma)] \\
            \absstatement{s_1; s_2}(S) = \absstatement{s_2}(\absstatement{s_1}(S)) \\
            \absstatement{\ite{c}{s_1}{s_2}}(S) =\\ \qquad\qquad\qquad \absstatement{s_1}(\abscexpression{c}(S)) \cup  \absstatement{s_1}(\abscexpression{\lnot c}(S)) \\
            \absstatement{\while{c}{s}}(S) =\\ \qquad\qquad\qquad \abscexpression{\lnot c}(\lfp (X \mapsto X \cup \absstatement{s}(\abscexpression{c}(X))))
        \end{array}
    \]
    }
    \uncover<4->{
        Wow, this is almost all computable now!!
        The only thing that isn't is the while loop...
        It's a join of infinitely many things!
    }
\end{frame}

\begin{frame}
    \frametitle{Fixpoint acceleration}
    If we have a sequence of intervals \(([0, 0], [0, 1], [0, 2], [0, 3], \dots)\), then there's no hope to find a fixed-point of them by joining them repeatedly.
    
    \uncover<2->{
    We can widen intervals like so:
    \[
        [a, b] \mathbin{\nabla} [c, d] = \left[\begin{cases}
            a & \text{if } a \leq c \\
            -\infty & \text{otherwise} \\
        \end{cases}, \begin{cases}
            b & \text{if } b \geq d \\
            \infty & \text{otherwise} \\
        \end{cases}\right].
    \]
    Intuitively, if we see \([c, d]\) is growing, clamp to infinity.

    This is an instance of what is called a widening operator (definition omitted)!
    }

    \uncover<3->{
    \begin{theorem}
        If \(g : A \to A\) is an approximation of \(f : C \to C\) (in the sense of the (lax) commutative diagram from earlier), then the sequence \(x_0 = \bot\) and \(x_{i+1} = x_i \mathbin{\nabla} g(x_i)\) stabilises in finite time to \(x\) to approximate the least fixed point of \(f\), \(\lfp(f) \sqsubseteq \gamma(x)\).
    \end{theorem}
    no proof
    }
\end{frame}

\begin{frame}
    \frametitle{Fixpoint acceleration}

    \uncover<2->{
    We can then define widening on \(I^V\) pointwise, and define while semantics with
    \[
        \begin{array}{l}
            \absstatement{\while{c}{s}}(S) =\\ \qquad\qquad\qquad \abscexpression{\lnot c}(\lfp (X \mapsto X \mathbin{\nabla} (S \sqcup \absstatement{s}(\abscexpression{c}(X))))).
        \end{array}
    \]
    This is computable!
    }

    \uncover<3->{
    Also everything here is an approximation of the concrete semantics i promise
    }
\end{frame}

\begin{frame}
    \frametitle{Abstract interpreting}
    \only<-20>{
    \begin{algorithmic}
        \State{\(F \gets \underbrace{[0, \dots, 0]}_{\text{100 of them}}\) \uncover<1->{\textcolor{gray}{\(i \mapsto [\infty, \infty]\)}}}
        \State{\(F[0] \gets 0\) \uncover<2->{\textcolor{gray}{\(i \mapsto [\infty, \infty]\)}}}
        \State{\(F[1] \gets 1\) \uncover<3->{\textcolor{gray}{\(i \mapsto [\infty, \infty]\)}}}
        \State{\(i \gets 2\) \uncover<4->{\textcolor{gray}{\(i \mapsto [2, 2]\)}}}
        \While{\(i < 100\)}
            \State{\(
                \only<-5>{\uncover<5-5>{\textcolor{gray}{i \mapsto [2, 2] \sqcap [-\infty, 99]}}}
                \only<6-11>{\uncover<6-11>{\textcolor{gray}{i \mapsto [2, 2]}}}
                \only<12-12>{\uncover<12-12>{\textcolor{gray}{i \mapsto [2, \infty] \sqcap [-\infty, 99]}}}
                \only<13->{\uncover<13->{\textcolor{gray}{i \mapsto [2, 99]}}}
            \)}
            \only<-19>{
            \If{\(2 > i \lor i > 99\)}
                \State \textbf{explode} (out of bounds)
            \EndIf
            }
            \State{\(F[i] \gets F[i-1] + F[i-2]\) 
            \(
                \only<-13>{\uncover<7-13>{\textcolor{gray}{i \mapsto [2, 2]}}}
                \only<14->{\uncover<14->{\textcolor{gray}{i \mapsto [2, 99]}}}
            \)}
            \State{\(i \gets i + 1\)
            \(
                \only<-14>{\uncover<8->{\textcolor{gray}{i \mapsto [3, 3]}}}
                \only<15->{\uncover<15->{\textcolor{gray}{i \mapsto [3, 100]}}}
            \)}
            \State{\(
                \only<1-9>{\uncover<9-9>{\textcolor{gray}{i \mapsto [2, 2] \mathbin{\nabla} ([2, 2] \sqcup [3, 3])}}}
                \only<10-10>{\uncover<10-10>{\textcolor{gray}{i \mapsto [2, 2] \mathbin{\nabla} [2, 3]}}}
                \only<11-15>{\uncover<11-15>{\textcolor{gray}{i \mapsto [2, \infty]}}}
                \only<16-16>{\uncover<16-16>{\textcolor{gray}{i \mapsto [2, 99] \mathbin{\nabla} ([2, 2] \sqcup [3, 100])}}}
                \only<17-17>{\uncover<17-17>{\textcolor{gray}{i \mapsto [2, 99] \mathbin{\nabla} [2, 100]}}}
                \only<18->{\uncover<18->{\textcolor{gray}{i \mapsto [2, \infty]}}}
            \)}
        \EndWhile
        \State{\(\uncover<19->{\textcolor{gray}{i \mapsto [100, \infty]}}\)}
        \\\noindent\Return \(F\)
    \end{algorithmic}
    }
    \only<21->{
    \begin{algorithmic}
        \State{\(F \gets \underbrace{[0, \dots, 0]}_{\text{100 of them}}\)}
        \State{\(F[0] \gets 0\)}
        \State{\(F[1] \gets 1\)}
        \State{\(i \gets 2\)}
        \While{\(i < 100\)}
            \State{\(F[i] \gets F[i-1] + F[i-2]\)}
            \State{\(i \gets i + 1\)}
        \EndWhile
        \\\noindent\Return \(F\)
    \end{algorithmic}
    we did it!!
    }
\end{frame}

\begin{frame}
    \frametitle{Further reading}
    \nocite{*}
    \bibliographystyle{plain}
    \bibliography{references}
\end{frame}

\end{document}
